\section{Purpose Definition}

This section reports the first iTelos phase. Its input is the informal purpose of Section 2 together with the three data sources named in the project proposal; its output is the formal purpose: scenarios, personas, competency questions (CQs), the concepts that will define entity types and properties, and the Purpose ER model. The phase was carried out \emph{middle-out}: descriptive attributes were derived from the CQs, while identifying attributes were derived from the actual columns of the datasets, which were inspected during information gathering (Section~\ref{sec:info-gathering}). All design material is stored in the repository folder \texttt{Phase 1 - Purpose Definition/}.

\subsection{Purpose Formalization}

\subsubsection{Scenarios}

Six usage scenarios describe how the KG is used after a spatial transcriptomics analysis. They are illustrative design artefacts: they do not describe real studies or real results, and gene, pathway, and tissue names are used as plausible examples (all of them were checked to exist in the source files).

\begin{table}[htbp]
\centering
\small
\begin{tabularx}{\linewidth}{|l|l|X|l|}
\hline
\textbf{ID} & \textbf{Persona} & \textbf{Scenario} & \textbf{CQs} \\ \hline
S1 & P1 & \textbf{Reading a tissue frontier.} At the frontier between a tumour cluster and the surrounding stroma, \emph{Epithelial--mesenchymal transition} and \emph{Hypoxia} score high on the tumour side; the analyst wants to see the genes of the two pathways before interpreting them. & CQ1 \\ \hline
S2 & P3 & \textbf{Learning the Hallmark collection.} A student wants the list of the 50 pathways with their genes and official names, to understand what a pathway score measures. & CQ1, CQ1.2 \\ \hline
S3 & P2 & \textbf{From a pathway to candidate cell types.} \emph{Interferon gamma response} increases along an invasive border; the immunologist wants the cell types whose markers overlap with the pathway genes, ideally limited to markers reported in the relevant tissue, to choose which populations to stain for. & CQ3, CQ3.1, CQ3.3 \\ \hline
S4 & P1 & \textbf{Activated versus suppressed pathways.} Across a frontier there are three activated and two suppressed pathways; the analyst wants the genes that belong to at least one pathway of each group. & CQ4, CQ4.1, CQ4.2 \\ \hline
S5 & P3 & \textbf{A single gene, many pathways.} \emph{MYC} appears among the differentially expressed genes; the student wants every Hallmark pathway that contains it. & CQ2, CQ2.1 \\ \hline
S6 & P4 & \textbf{Checking a panel gene.} In an image-based panel, \emph{CD74} lights up in one region; the pathologist wants its approved HGNC identity, the cell types it marks in that tissue, and its pathways. & CQ2, CQ2.2 \\ \hline
\end{tabularx}
\caption{Usage scenarios}
\label{tab:scenarios}
\end{table}

\subsubsection{Personas}

Four semi-fictional personas represent groups of end users (Table~\ref{tab:personas}). They are illustrative profiles following the iTelos persona technique and do not describe real people.

\begin{table}[htbp]
\centering
\small
\begin{tabularx}{\linewidth}{|l|l|X|X|}
\hline
\textbf{ID} & \textbf{Persona} & \textbf{Profile} & \textbf{Purpose} \\ \hline
P1 & Dr.\ Laura Ferri, 34 & Spatial transcriptomics bioinformatician in a core facility; her pipeline scores the 50 Hallmark pathways per region and returns activated and suppressed pathway lists. Uses Python/R and SPARQL. & Read the genes behind each pathway at a frontier; find genes shared by activated and suppressed pathways. \\ \hline
P2 & Dr.\ Marco Ferretti, 41 & Tumour immunologist (wet lab); receives results from bioinformaticians and plans follow-up experiments. Uses a web query interface; does not write code. & Find cell types whose marker genes overlap with a pathway that changes across a tumour border, as a hint, not as evidence of activity. \\ \hline
P3 & Giulia Neri, 24 & MSc student in bioinformatics, new to gene-set analysis, working on a public spatial dataset. & Browse the content of the Hallmark pathways; find the pathways that contain a given gene. \\ \hline
P4 & Dr.\ Chiara Bassi, 47 & Molecular pathologist reviewing targeted, image-based spatial panels gene by gene. & Confirm the official identity of a gene, the cell types it marks in a tissue, and its pathways. \\ \hline
\end{tabularx}
\caption{Personas (illustrative)}
\label{tab:personas}
\end{table}

\subsubsection{Competency Questions}

CQ1--CQ4 come from the project proposal. To make the requirement list more varied, nine sub-questions refining them were proposed; seven were approved and two were discarded (Section~\ref{sec:qual-eval}). Sub-question identifiers were kept unchanged after the discard, for traceability. All CQs use only the three entity types and the two relations of the scope.

\begin{table}[htbp]
\centering
\small
\begin{tabularx}{\linewidth}{|l|X|l|l|}
\hline
\textbf{ID} & \textbf{Competency question} & \textbf{Personas} & \textbf{Scenarios} \\ \hline
CQ1 & Which human genes are associated with each of the 50 MSigDB Hallmark pathways? & P1, P3 & S1, S2 \\ \hline
CQ1.2 & What are the HGNC ID and the approved name of each gene in a given pathway? & P3 & S2 \\ \hline
CQ2 & Which biological pathways contain a given human gene? & P3, P4 & S5, S6 \\ \hline
CQ2.1 & Which genes belong to more than one Hallmark pathway, and to how many? & P3 & S5 \\ \hline
CQ2.2 & For which cell types is a given gene reported as a marker, optionally within a given tissue class? & P4 & S6 \\ \hline
CQ3 & Which cell types have marker genes in common with the genes of a selected pathway? & P2 & S3 \\ \hline
CQ3.1 & Which cell types share marker genes with a selected pathway, considering only markers reported in a given tissue class? & P2 & S3 \\ \hline
CQ3.3 & Which specific genes make up the overlap between a given pathway and a given cell type? & P2 & S3 \\ \hline
CQ4 & Which genes are shared between pathways identified as activated and pathways identified as suppressed in an upstream spatial transcriptomics analysis? & P1 & S4 \\ \hline
CQ4.1 & For each gene shared by the two groups, which activated and which suppressed pathways contain it? & P1 & S4 \\ \hline
CQ4.2 & Which of the genes shared by the two groups are also markers of at least one cell type, and of which ones? & P1 & S4 \\ \hline
\end{tabularx}
\caption{Final competency questions}
\label{tab:cqs}
\end{table}

\noindent Two semantic constraints apply. The answer to CQ3 is a descriptive gene overlap, not evidence of pathway activity in the cell type. CQ4 takes the activated and suppressed pathway lists as query inputs: a gene answers CQ4 if it belongs to at least one activated \emph{and} at least one suppressed pathway; the KG does not decide which pathways are activated or suppressed.

\subsubsection{Informal labels and popularity categories}

The concepts were collected in a Purpose Formalization sheet (\texttt{pf\_sheet.csv}), which links each concept to the scenarios, personas, and CQs that require it. Each concept was classified with the Focus parameter:
\begin{itemize}
    \item \textbf{Core}: \emph{Pathway} and \emph{CellType}, and the two relations. The purpose is built around linking pathways and cell types through genes.
    \item \textbf{Common}: \emph{Gene}. It is a general entity, widely modelled by reference resources (HGNC), and it is the bridge between pathways and cell types.
    \item \textbf{Contextual}: the descriptive context of cell types (specific disease, tissue type, broader cell class) and the evidence attached to marker links.
\end{itemize}

\begin{table}[htbp]
\centering
\small
\begin{tabularx}{\linewidth}{|l|l|X|X|}
\hline
\textbf{EType} & \textbf{Focus} & \textbf{Identifying attributes (from datasets)} & \textbf{Descriptive attributes (from CQs)} \\ \hline
Gene & Common & HGNC ID (e.g.\ \texttt{HGNC:1697}); approved symbol as secondary & approved symbol, approved name (CQ1.2); source symbol(s) and resolution status (Contextual, audit) \\ \hline
Pathway & Core & Hallmark gene-set identifier (e.g.\ \texttt{HALLMARK\_HYPOXIA}); MSigDB URL as secondary & label (e.g.\ \emph{Hypoxia}), source URL \\ \hline
CellType & Core & (cell name, tissue class, disease category), e.g.\ \emph{Macrophage -- Lung -- Cancer}; Cell Ontology ID(s) as secondary & cell name, tissue class (CQ2.2, CQ3.1), disease category; disease, tissue type, cell name class (Contextual, multi-valued where needed); Cell Ontology ID (Common) \\ \hline
\end{tabularx}
\caption{Entity types, Focus, and attributes}
\label{tab:etypes}
\end{table}

\begin{table}[htbp]
\centering
\small
\begin{tabularx}{\linewidth}{|l|l|l|X|}
\hline
\textbf{Relation} & \textbf{Domain $\rightarrow$ Range} & \textbf{Focus} & \textbf{Relation attributes} \\ \hline
participates in & Gene $\rightarrow$ Pathway (N:M) & Core & none (membership as stated by MSigDB) \\ \hline
is marker for & Gene $\rightarrow$ CellType (N:M) & Core & marker source, method details, PubMed ID, publication year (Contextual) \\ \hline
\end{tabularx}
\caption{Relations}
\label{tab:relations}
\end{table}

\noindent The main design choices, with their strong and weak points, were the following.
\begin{itemize}
    \item \textbf{Cell type naming.} A cell type is named by the CellMarker cell name rather than by its Cell Ontology (CL) identifier. Names and CL identifiers are not one-to-one in the source (\emph{Macrophage} occurs with 24 different CL identifiers), and the name is what users search for. The CL identifiers are kept as a multi-valued attribute for linking to the Cell Ontology. Weak point: two names for the same biological concept remain separate.
    \item \textbf{Context-qualified cell types.} The same cell type is reported in different tissues and disease contexts (\emph{Macrophage} appears in 98 tissue classes). Pooling all of them into one context-free marker profile would lose this distinction and would prevent the tissue filter of CQ2.2 and CQ3.1. The identity of a CellType is therefore the triple (cell name, tissue class, disease category), which gives 12,632 CellTypes; with the cell name alone there would be 3,190. The three parts remain separate literals, so all CellTypes with a given base name can still be retrieved. Tissue class and disease category were chosen because they are always present and have few, clean values (138 tissue classes; disease category is one of \emph{Normal}, \emph{Cancer}, \emph{Disease}); no mapping to external tissue or disease ontologies is needed. Weak point: the number of CellTypes grows about fourfold.
    \item \textbf{Specific disease as a descriptive attribute.} The specific disease is free text with 1,067 values (e.g.\ \emph{Lung adenocarcinoma (LUAD)}); using it in the identity would require normalization. It is kept as a multi-valued attribute: 2,126 of the 12,632 CellTypes have more than one disease value (at most 93).
    \item \textbf{All marker evidence kept, with its source.} 91.9\% of the CellMarker records are computationally predicted markers; the rest are manually curated. All of them are kept, and the marker source is recorded on the \emph{is marker for} relation, so that queries can be restricted to curated markers. How to encode relation attributes in RDF without adding a fourth entity type is left to Phase 3.
\end{itemize}

\subsubsection{Purpose ER model}

Figure~\ref{fig:er} shows the Purpose ER model. Both relations are many-to-many: a gene can belong to several pathways (e.g.\ \emph{MYC} belongs to 10) and mark several cell types, and a pathway contains between 32 and 200 genes. Pathways and cell types are never linked directly: their connection (CQ3) is always computed through shared genes, so the model does not assert pathway activity in any cell type. The activated and suppressed pathway groups of CQ4 are not part of the model, since they are passed to the query.

\begin{figure}[htbp]
\centering
\begin{tikzpicture}[
    font=\scriptsize,
    entity/.style={rectangle, draw, thick, align=left, inner sep=3pt, text width=3.1cm},
    rel/.style={diamond, draw, thick, aspect=1.6, inner sep=0pt, align=center, fill=green!15},
    node distance=0.6cm and 0.3cm]
  \node[entity, fill=blue!12] (gene) {\textbf{GENE} \hfill \textit{Common}\\[2pt]
    $\ast$\,hgnc\_id\\ approved\_symbol\\ approved\_name\\ source\_symbols [multi]\\ resolution\_status};
  \node[rel, left=of gene] (part) {participates\\in};
  \node[entity, fill=green!15, left=of part, text width=2.6cm] (pathway) {\textbf{PATHWAY} \hfill \textit{Core}\\[2pt]
    $\ast$\,hallmark\_id\\ label\\ source\_url};
  \node[rel, right=of gene] (marker) {is marker\\for};
  \node[entity, fill=green!15, right=of marker, text width=3.4cm] (cell) {\textbf{CELLTYPE} \hfill \textit{Core}\\[2pt]
    $\ast$\,cell\_name\\ $\ast$\,tissue\_class\\ $\ast$\,disease\_category\\ disease [multi]\\ tissue\_type [multi]\\ cell\_name\_class\\ cell\_ontology\_id [multi]};
  \node[entity, fill=red!12, below=0.6cm of marker, text width=3.3cm] (ev) {\textit{relation attributes (Contextual)}\\ marker\_source, method\_details,\\ pmid [multi], year [multi]};
  \draw[thick] (gene) -- node[above]{N} (part) -- node[above]{M} (pathway);
  \draw[thick] (gene) -- node[above]{N} (marker) -- node[above]{M} (cell);
  \draw[dashed] (marker) -- (ev);
\end{tikzpicture}
\caption{Purpose ER model. Attributes marked with $\ast$ form the identity of the entity type; [multi] marks multi-valued attributes; colours show the Focus (blue: Common, green: Core, red: Contextual).}
\label{fig:er}
\end{figure}

\subsection{Information Gathering}
\label{sec:info-gathering}

\subsubsection{KG/KnowDive data sources}

\medskip\noindent\textbf{Knowledge layer.} No KG/KnowDive knowledge resource was collected during this phase. The selection of knowledge-layer resources is completed in Phases 2 and 3.

\medskip\noindent\textbf{Data layer.} No KG/KnowDive dataset was collected: all the data used by the project comes from the external sources below.

\subsubsection{External data sources}

\medskip\noindent\textbf{Knowledge layer.} The reference ontologies named in the project proposal (Gene Ontology, Cell Ontology, Pathway Ontology, Human Disease Ontology, NCI Thesaurus, through the OBO Foundry and NCBO BioPortal) were considered. Of these, only the Cell Ontology is directly relevant to the agreed scope, because CellMarker already provides CL identifiers for its cell types; it is kept as a \emph{Common} resource, used for the Cell Ontology ID attribute of CellType. Tissue and disease ontologies (UBERON, DOID) were not adopted, since tissue and disease are kept as literal context; the other ontologies concern concepts that are outside the scope (Section~2).

\medskip\noindent\textbf{Data layer.} The data layer consists of three external resources, downloaded on 2026-11-05 and inspected read-only on 2026-11-12. Their provenance and licences are recorded in the repository (\texttt{SOURCES.md}).

\begin{table}[htbp]
\centering
\small
\begin{tabularx}{\linewidth}{|>{\raggedright\arraybackslash}X|l|l|r|l|}
\hline
\textbf{Source (file)} & \textbf{Role} & \textbf{Format} & \textbf{Records} & \textbf{Licence} \\ \hline
MSigDB Hallmark \cite{liberzon2015molecular} (\texttt{h.all.v2026.2.Hs.symbols.gmt}) & content: Pathway, participates in & GMT & 50 gene sets & CC BY 4.0 \\ \hline
CellMarker 3.0 \cite{hu2026cellmarker} (\texttt{human\_cell\_marker.txt}) & content: CellType, is marker for & TSV, 23 columns & 1,779,944 & not stated \\ \hline
HGNC \cite{seal2026genenames} (\texttt{hgnc\_complete\_set.txt}) & reference: gene identity & TSV, 53 columns & 45,326 genes & CC0 \\ \hline
\end{tabularx}
\caption{External data sources}
\label{tab:sources}
\end{table}

\medskip\noindent\textbf{MSigDB Hallmark.} One gene set per line: name (always \texttt{HALLMARK\_...}), a URL to the gene-set page (there is no textual description), and the gene symbols. The 50 sets contain 32 to 200 genes, for 7,322 gene--pathway memberships and 4,384 distinct symbols, with no empty fields and no duplicates. 4,383 symbols are current HGNC approved symbols; one (\texttt{TARP}) matches no approved, previous, or alias symbol.

\medskip\noindent\textbf{HGNC.} All 45,326 records have unique HGNC IDs and symbols. The columns needed for identification (\texttt{hgnc\_id}, \texttt{symbol}, \texttt{name}, and, for diagnostics, \texttt{prev\_symbol} and \texttt{alias\_symbol}) are all present. Previous symbols and aliases are pipe-separated lists in double quotes. The other 44 columns are out of scope.

\medskip\noindent\textbf{CellMarker 3.0.} All records are human. Tissue class, tissue type, disease category, disease, cell name, and marker are never empty; the Cell Ontology identifier is missing in 0.2\% of the records, and the standardized gene symbol in 9.1\%. 91.9\% of the records come from computational marker-detection methods (e.g.\ COSG, CelliD, FindAllMarker), the rest from single-cell sequencing studies, experiments, reviews, and companies. Several data-quality issues were found and are left to Phase 2: 3,725 distinct symbols differ from an HGNC approved symbol only by letter case; 26,181 distinct markers without a standardized symbol match nothing in HGNC (e.g.\ clone names such as \texttt{RP11-1143G9.4}); a few CL identifiers use the colon form or contain UBERON identifiers; and 70,731 distinct record contents occur more than once.

\noindent Table~\ref{tab:key-columns} lists the columns of each source that the model uses; the other columns are not imported.

\begin{table}[htbp]
\centering
\small
\begin{tabularx}{\linewidth}{|l|>{\raggedright\arraybackslash}X|X|}
\hline
\textbf{Source} & \textbf{Column(s)} & \textbf{Used for} \\ \hline
MSigDB Hallmark & field 1 (gene-set name) & Pathway identifier, e.g.\ \texttt{HALLMARK\_HYPOXIA} \\ \hline
MSigDB Hallmark & field 2 (gene-set URL) & Pathway source URL \\ \hline
MSigDB Hallmark & fields 3+ (gene symbols) & \emph{participates in} (after HGNC resolution) \\ \hline
HGNC & \texttt{hgnc\_id}, \texttt{symbol}, \texttt{name} & Gene identity, approved symbol, approved name \\ \hline
HGNC & \texttt{prev\_symbol}, \texttt{alias\_symbol} & diagnostics for symbol resolution (Phase 2); not imported \\ \hline
CellMarker 3.0 & \texttt{cell\_name}, \texttt{tissue\_class}, \texttt{disease\_category} & CellType identity \\ \hline
CellMarker 3.0 & \texttt{disease}, \texttt{tissue\_type}, \texttt{cell\_name\_class}, \texttt{cellontology\_id} & CellType descriptive attributes \\ \hline
CellMarker 3.0 & \texttt{symbol} (or \texttt{marker} when \texttt{symbol} is empty) & gene of \emph{is marker for} (after HGNC resolution) \\ \hline
CellMarker 3.0 & \texttt{marker\_source}, \texttt{method\_details}, \texttt{pmid}, \texttt{year} & attributes of \emph{is marker for} \\ \hline
\end{tabularx}
\caption{Source columns used by the Purpose ER model}
\label{tab:key-columns}
\end{table}

\medskip\noindent\textbf{Classification.} MSigDB Hallmark is the source of the Core \emph{Pathway} type and its relation; CellMarker is the source of the Core \emph{CellType} type, its Contextual attributes, and the \emph{is marker for} relation; HGNC is the source of the Common \emph{Gene} identity. The collected resources cover all approved CQs (Table~\ref{tab:coverage}), so no further source identification was needed.

\begin{table}[htbp]
\centering
\small
\begin{tabularx}{\linewidth}{|l|X|X|}
\hline
\textbf{CQs} & \textbf{Information needed} & \textbf{Source} \\ \hline
CQ1, CQ1.2 & pathway $\rightarrow$ genes; HGNC ID and approved name & MSigDB; HGNC \\ \hline
CQ2, CQ2.1 & gene $\rightarrow$ pathways & MSigDB \\ \hline
CQ2.2, CQ3.1 & gene $\rightarrow$ cell types; tissue filter & CellMarker (cell name, tissue class) \\ \hline
CQ3, CQ3.3 & pathway genes $\cap$ cell-type marker genes & MSigDB and CellMarker, joined through HGNC \\ \hline
CQ4, CQ4.1, CQ4.2 & pathway groups (query input) $\rightarrow$ shared genes and their markers & MSigDB (and CellMarker) \\ \hline
\end{tabularx}
\caption{Coverage of the competency questions by the collected sources}
\label{tab:coverage}
\end{table}

\subsection{Qualitative Evaluation}
\label{sec:qual-eval}

\medskip\noindent\textbf{Scenarios.} Six scenarios were initially considered and six were finally kept. One was revised: the example gene of S6 was changed from \emph{CD68} to \emph{CD74}, because \emph{CD68} is a CellMarker marker but belongs to no Hallmark pathway, so the pathway part of the scenario would have had an empty answer.

\begin{table}[htbp]
\centering
\small
\begin{tabularx}{\linewidth}{|l|X|l|}
\hline
\textbf{ID} & \textbf{Scenario} & \textbf{Status} \\ \hline
S1 & Reading a tissue frontier & Retained \\ \hline
S2 & Learning the Hallmark collection & Retained \\ \hline
S3 & From a pathway to candidate cell types & Retained \\ \hline
S4 & Activated versus suppressed pathways & Retained \\ \hline
S5 & A single gene, many pathways & Retained \\ \hline
S6 & Checking a panel gene & Revised (example gene \emph{CD68} $\rightarrow$ \emph{CD74}) \\ \hline
\end{tabularx}
\caption{Scenarios: initial vs final (6 initial, 6 final)}
\end{table}

\medskip\noindent\textbf{Personas.} Four personas were initially considered and four were finally kept. 

\begin{table}[htbp]
\centering
\small
\begin{tabularx}{\linewidth}{|l|X|l|}
\hline
\textbf{ID} & \textbf{Persona} & \textbf{Status} \\ \hline
P1 & Dr.\ Laura Ferri, spatial transcriptomics bioinformatician & Retained \\ \hline
P2 & Dr.\ Marco Ferretti, tumour immunologist & Retained \\ \hline
P3 & Giulia Neri, MSc student in bioinformatics & Retained \\ \hline
P4 & Dr.\ Chiara Bassi, molecular pathologist & Retained \\ \hline
\end{tabularx}
\caption{Personas: initial vs final (4 initial, 4 final)}
\end{table}

\medskip\noindent\textbf{Competency questions.} Thirteen CQs were initially considered: the four of the project proposal and nine proposed sub-questions. Eleven were finally kept. CQ1.1 (number of genes per pathway) and CQ3.2 (number of shared genes per overlapping cell type) were discarded because they were deemed not relevant to the project's goals.

\begin{table}[htbp]
\centering
\small
\begin{tabularx}{\linewidth}{|l|X|l|}
\hline
\textbf{ID} & \textbf{Competency question (short)} & \textbf{Status} \\ \hline
CQ1 & Genes of each Hallmark pathway & Retained \\ \hline
CQ1.1 & Number of genes per pathway & Discarded \\ \hline
CQ1.2 & HGNC ID and approved name of the genes of a pathway & Retained \\ \hline
CQ2 & Pathways containing a gene & Retained \\ \hline
CQ2.1 & Genes in more than one pathway & Retained \\ \hline
CQ2.2 & Cell types marked by a gene (optionally per tissue class) & Retained \\ \hline
CQ3 & Cell types sharing marker genes with a pathway & Retained \\ \hline
CQ3.1 & CQ3 restricted to a tissue class & Retained \\ \hline
CQ3.2 & Number of shared genes per overlapping cell type & Discarded \\ \hline
CQ3.3 & Genes forming a pathway--cell-type overlap & Retained \\ \hline
CQ4 & Genes shared by activated and suppressed pathways & Retained \\ \hline
CQ4.1 & Pathways of each group containing a shared gene & Retained \\ \hline
CQ4.2 & Shared genes that are also cell-type markers & Retained \\ \hline
\end{tabularx}
\caption{Competency questions: initial vs final (13 initial, 11 final)}
\end{table}

\medskip\noindent\textbf{Purpose refinement.} The informal purpose did not change in substance. The temporal boundary was made explicit (1978--2026) from the publication years in CellMarker, and the ER model was adjusted after information gathering: the inspection of the data showed that cell names and CL identifiers are not one-to-one and that the same cell type occurs in many tissue and disease contexts, which led to the context-qualified \emph{CellType} described above.

\medskip\noindent\textbf{Dataset-level formatting and transformations.} None. In this phase the raw files were only inspected, in read-only mode; no file was modified, filtered, or converted. The parsing rules proposed for Phase 2 are: select columns by name; strip the double quotes of HGNC multi-valued fields before splitting them; use the CellMarker standardized symbol when present and the original marker otherwise, keeping the original string for audit; normalize CL identifiers of the form \texttt{CL:nnnnnnn} to \texttt{CL\_nnnnnnn} and treat non-CL identifiers as missing; and never drop records silently.

\medskip\noindent\textbf{Remarks on LLM usage.} LLMs were used often in this phase (Appendix A). The Claude Code agent was used for brainstorming, for refining the Markdown design files in the phase folder, and for writing this documentation; Gemini was used earlier, while writing the project proposal, to identify possible data sources. Among the phase tasks, LLM support was occasional for generating scenarios and personas, for the qualitative evaluation of the purpose and of the data sources, for identifying data sources, and for resolving schema ambiguities. Every textual output was checked by hand, and numerical outputs were checked at least for consistency. Few errors were found, mostly errors of context; they were addressed at the root by providing the agent with the course guides and with a project instruction file (\texttt{AGENT.md}).

\noindent\textbf{The phase-specific LLM questionnaire is answered in Appendix A.}
